% ============================================================
%  Chapter: Consequences of Ethics, Applied Ethics
%  Foundation Batch --- Lectures 3 \& 4
%  Topics: Consequences of Ethics \textbullet\ Applied Ethics \textbullet\ Dimensions of Ethics
%          \textbullet\ Ethical Relativism \& Absolutism \textbullet\ Environmental Ethics
%          \textbullet\ Technology Ethics \textbullet\ Media Ethics \textbullet\ Public vs Private Relationship
% ============================================================

\chapter{Consequences of Ethics and Applied Ethics}

\vspace{14pt}

\section{Consequences of Ethics}

Consequences are one of the most effective ways to make people understand the importance of any particular thing --- including ethics. Warning labels on cigarette packs, a teacher warning students that they will fail if they do not practise answer writing, religious traditions warning of divine punishment --- all these are appeals to consequences. The same logic applies to ethics: why should one be ethical? Because of the consequences of being ethical (and of not being ethical).

Consequences operate at multiple levels: individual, organisational, societal, and national/global.

\subsection{Why Individuals Should Be Ethical}

\paragraph{Survival} Being ethical is important for our very biological survival. Formal institutions of the state --- police, judiciary, and the military --- will eliminate those who are grossly unethical. Terrorists and extremists are killed because they have crossed the minimum threshold of ethical acceptability for society. Drug and alcohol addiction destroy physical health; unethical use of one's own body is, therefore, ultimately harmful to one's survival.

\begin{ExampleBox}
Drug addiction and alcohol addiction impair physical health --- a person who is addicted is not being dutiful towards himself. If you indulge in heinous crimes, society will not allow you to survive; formal institutions will eliminate you.
\end{ExampleBox}

\paragraph{Psychological Well-Being} Even when unethical behaviour goes unpunished externally, it exacts a psychological price. Conscience generates guilt, regret and shame. A person who has committed corruption, for instance, may profit materially, but deep down they know what they have done --- and long-term psychological peace is impossible without ethical conduct. Choosing not to discharge one's duty may feel comfortable in the short term, but will not produce happiness in the long run.

\begin{KeyConcept}
Psychological well-being requires that we examine our conduct. An unexamined life eventually produces dissatisfaction even if material success is achieved.
\end{KeyConcept}

\paragraph{Quality of Interpersonal Relationships} Human beings are social animals; the quality of our relationships is a major determinant of the quality of our life. Ethics --- especially honesty and trustworthiness --- is the foundation of good relationships. The classic illustration is the story of ``The Boy Who Cried Wolf'': habitual lying destroys credibility, and when the wolf actually comes, no one trusts him. The consequence is isolation and helplessness.

\begin{DataFact}
\begin{itemize}
\item Rising divorce rates and the breakdown of family structures are partly consequences of the erosion of trust in intimate relationships.
\item Trust breakdown in relationships, sometimes labelled as the ``Blue Drum Syndrome'', is a consequence of declining ethical conduct in personal life.
\end{itemize}
\end{DataFact}

\paragraph{Professional Success and Fulfilment} Ethics enables professional success in two distinct ways: instrumental success (promotion, credibility, leadership) and intrinsic fulfilment (deriving meaning and satisfaction from work). When you are honest, duty-bound, and committed, your profession becomes a source of meaning.

\begin{ExampleBox}
\begin{itemize}
\item ``Missile Man'' Dr A.P.J. Abdul Kalam --- the depth of his commitment to his duty made science his greatest source of fulfilment; his career and character became inseparable.
\item A ranker from 2013 (Pawan Singh Rana, rank 249) who came from a very humble background in Bihar; when such students succeed, it is the outcome of commitment to purpose and duty. The satisfaction on seeing a student succeed is itself an example of how a profession becomes a source of fulfilment when one is duty-bound.
\end{itemize}
\end{ExampleBox}

Ethics is connected to professional growth in concrete economic terms as well: ethical conduct builds the credibility needed to receive better opportunities, command respect, attract cheaper credit, and earn trust that leads to leadership.

\begin{KeyConcept}
Nishkaam Karma (Bhagavad Gita): do your duty without attachment to the fruits. When work is done with this spirit, it becomes a source of meaning rather than merely a means to an end.
\end{KeyConcept}

\paragraph{Leadership and Credibility} Respect cannot be commanded; it must be earned. Great leaders --- Mahatma Gandhi, Swami Vivekananda, Ramakrishna Paramahamsa, Subhash Chandra Bose, Bhagat Singh --- earned respect through ethical conduct. This is true at every level: a manager who is committed and ethical over years becomes the person the organization can rely on, the one who gets promoted, the one whose team follows with genuine commitment.

\paragraph{Living a Purposeful Life --- Going Beyond Animalistic Existence} Human beings are higher-order animals. A tiger that kills a human is never judged as wrong; but a human being who kills another is judged --- because human beings are endowed with higher cognitive faculties that allow us to distinguish right from wrong. Ethics calls on us to use those faculties; it lifts us beyond mere animalistic existence to a life with purpose.

\begin{PYQLink}
\begin{itemize}
\item ``Life is a long journey between human being and being human.'' --- UPSC topic. All of us are born human beings; to live a life worthy of being called human, we must be ethical.
\item ``It is better to be a human being dissatisfied than a pig satisfied.'' (J.S. Mill.) A pig cannot distinguish between higher-quality and lower-quality pleasures; a human being, by contrast, has an ethical compass that enables the choice of higher pleasures even at the cost of dissatisfaction. A rapist, a cheater, a paper-leaker behaves more like a pig --- gratifying a lower impulse.
\item ``The unexamined life is not worth living.'' (Socrates.) Ethics compels introspection; introspection makes life purposeful.
\end{itemize}
\end{PYQLink}

\begin{ExampleBox}
\begin{itemize}
\item Valmiki (Ramayana): once a dacoit (Ratnakar), he was robbing people to feed his family. Naradmuni asked him whether his family would share in the punishment for the crimes he committed to feed them. They refused. This made him realise the consequences of his unexamined life, and he transformed. The same introspective journey changed his life completely.
\item King Ashoka: after the Kalinga War, introspection transformed him into ``Ashoka the Great'' --- he propagated peace through the Ashokan Pillars.
\item Prince Siddhartha became Gautam Buddha through the same process: examining the reality of suffering led to transformation.
\item Narendra Nath Datta became Swami Vivekananda.
\item Meer Jafar, Dola Rai --- examples of those who chose betrayal; history's memory of them illustrates the legacy consequences of unethical choices.
\end{itemize}
\end{ExampleBox}

\begin{KeyConcept}
\begin{itemize}
\item Ethics helps us in realising our potential --- what Socrates called Eudaimonia: the fulfilment of higher-order needs. As a civil services aspirant, you will never realise your potential if you do not examine and improve your conduct.
\item Ethics allows us to create a legacy that outlasts our lives.
\end{itemize}
\end{KeyConcept}

\begin{QuickRevision}
Individuals should be ethical for: (1) survival, (2) physical well-being, (3) psychological well-being, (4) quality of interpersonal relationships, (5) professional success and credibility, (6) leadership, (7) purposeful life / Eudaimonia, (8) creating legacy.
\end{QuickRevision}

\subsection{Why Organisations Should Be Ethical}

The same reasoning extends to organisations, both private and public. The consequences of organisational ethics (and its absence) operate in several dimensions.

\begin{itemize}
\item Brand value and consumer trust: an unethical organisation loses brand value and the trust of consumers. Without trust, word-of-mouth collapses, marketing costs rise, and revenue falls.
\item Work culture, talent retention and motivation: discrimination on the basis of caste, religion or gender in an organisation destroys morale and motivation. People will not want to work hard if meritocracy is absent. Conversely, an organisation that promotes employee well-being, mental health, and innovation --- for example, where thoughtful failures are not punished --- attracts the best talent and drives productivity.
\item Efficiency and sustainability: corruption within an organisation creates severe inefficiency. As Rajiv Gandhi noted, when 1 rupee is spent, only 16 paise reaches the intended beneficiary. Blue Smart cab case --- the founder misappropriated investor funds to buy a personal flat worth \rupee 100 crore; the company is now shutting down. Satyam scam --- the company colluded with its auditors to falsify accounts and show inflated profits; the deception unravelled.
\end{itemize}

\begin{CriticalPoint}
``You can fool some of the people some of the time, but not all the people all the time.'' Unethical behaviour in organisations ultimately surfaces and causes catastrophic collapse.
\end{CriticalPoint}

\begin{itemize}
\item Good governance tools: JAM trinity (Jan Dhan, Aadhaar, Mobile) and the Government e-Marketplace (GeM) portal for tendering are examples of technology deployed to make procurement transparent and curb corruption --- ensuring fair price discovery.
\item Accountability and role modelling: when employees are held accountable and leadership models ethical behaviour, motivation is high, there is strong teamwork, and the organisation commands credibility.
\end{itemize}

\subsection{Why Society Should Be Ethical}

\begin{itemize}
\item Tolerance: tolerance means respecting difference, even if you do not accept it. Without tolerance, science and art cannot grow --- Galileo's heliocentric theory was condemned because of religious intolerance for a different scientific opinion. Censorship and intolerance stifle scientific temper, innovation, and artistic expression. Mob lynching and communal violence are direct consequences of the absence of tolerance in society.
\item Social harmony and fraternity: harmony means that different communities, castes, and groups can co-exist. Without it, divisions grow, leading to conflict and collective harm.
\item Social mobility: in an ascriptive society --- where what you achieve is decided by the identity you are born with (caste, gender, religion) --- mobility is very low. No matter how talented a person is, they cannot move out of the occupation assigned by birth. Ethics demands fairness and equal opportunity; without it, talented people are wasted.
\end{itemize}

\begin{ExampleBox}
Women: many talented women quit their careers because they bear a disproportionate share of household work --- physical mobility itself is restricted for women in many contexts. If women cricketers face structural disadvantages from childhood (training camps full of boys, poor infrastructure), the resulting performance gap cannot be attributed simply to lower ability. We must look beyond the objective measure to the structural inequity behind it. If 100 men and 100 women have equal average IQ, but men are given 10 times more opportunity to develop skills, the outcome will falsely appear meritocratic.
\end{ExampleBox}

\begin{itemize}
\item Rule of law: mob lynching is what happens when people lose faith in rule of law and take justice into their own hands. If people are not ethical enough to believe in and follow institutions, rule of law collapses.
\item Innovation: intolerance and censorship suppress thought, which suppresses innovation.
\item Cultural foundation: art and culture grow only in an environment of freedom and tolerance. Censorship damages cultural production.
\item Sustainable development: Bhutan collectively measures Gross National Happiness rather than only GDP --- a societal ethical choice to value well-being over mere growth.
\item Public health: open defecation, poor waste management, air pollution, and malnutrition in women due to patriarchal food-distribution norms are all public-health consequences of societal ethical failures.
\item Social capital: social capital is the network of relationships and mutual trust in a society. Where social capital is high --- as in a village community where the entire neighbourhood participates in a wedding, or the custom of Kanyadaan where even uninvited people contribute --- the community can collectively face adversity. Distrust and individualism erode social capital; trust-based societies need fewer formal enforcement mechanisms, have lower crime, and show greater co-operative capacity.
\item Trust in institutions: when government and public institutions are unethical, they lose legitimacy. Citizens withdraw trust, common goods are not provided, and the state loses its reason for existence.
\end{itemize}

\begin{QuickRevision}
\begin{itemize}
\item Societal ethics produces: tolerance, social harmony, social mobility, rule of law, innovation, cultural growth, sustainability, public health, social capital, and trust in institutions.
\item Conversely, unethical societies exhibit mob lynching, communal violence, gender and caste discrimination, corruption, and declining social capital.
\end{itemize}
\end{QuickRevision}

\section{Dimensions of Ethics}

Ethics has both a theoretical dimension and an applied/practical dimension. The theoretical dimension deals with the foundational philosophical standards and frameworks; the applied dimension is the application of those ethical principles in specific professional or social domains.

\subsection{Normative (Theoretical) Dimension}

A norm is a prescription of behaviour --- a guideline about what one is supposed to do. Normative ethics is therefore prescriptive ethics: it tells us what ethical standards we should follow. Just as different family members give different advice about what to do, different schools of thought within normative ethics prescribe different standards.

\begin{KeyConcept}
The three major normative frameworks are: (1) Deontological ethics, (2) Teleological ethics, and (3) Virtue ethics.
\end{KeyConcept}

\paragraph{Deontological Ethics} Focus: means and duty. An action is ethical if the means used are right and if one is doing one's duty --- regardless of the consequences. The rightness of an action lies in the action itself, not in its outcomes.

\begin{ExampleBox}
\begin{itemize}
\item If you cheat in an exam to clear it, deontology says the action is wrong --- because the means (cheating) are impure, even though the outcome (clearing the exam) may be desirable.
\item Poonam Pandey case --- she faked her own death to raise awareness about cervical cancer. A deontologist says: the action was wrong because she lied. Means matter more than ends; you cannot do the right thing through the wrong means.
\end{itemize}
\end{ExampleBox}

\paragraph{Teleological Ethics (Consequentialism)} Focus: consequences and ends. An action is ethical if its consequences are good. Means are secondary; what matters is whether the outcome is right.

\begin{ExampleBox}
\begin{itemize}
\item Poonam Pandey case --- a teleologist says: the action was right, because it raised awareness about cervical cancer, no one died, and the ends (public health awareness) were good.
\item ``I didn't know my food had sugar in it'' --- teleological ethics says you should be mindful of consequences; ignorance of consequences is not an excuse because you are responsible for them.
\end{itemize}
\end{ExampleBox}

\paragraph{Virtue Ethics} Focus: character of the individual. An action is ethical if it is performed by a person of good character, and if that action contributes to the development of good character. It asks not ``did the action produce the right outcome?'' or ``were the means right?'' but ``what kind of person are you becoming through this action?''

\begin{KeyConcept}
\begin{itemize}
\item Virtues are positive character traits --- honesty, loyalty, compassion, discipline, courage. Vices are negative character traits --- greed, cowardice, dishonesty. Virtue ethics demands that we maximise our virtues and minimise our vices.
\item ``We are what we repeatedly do.'' Virtue ethics focuses on character formation through habit. Mind your habits, not just isolated actions.
\end{itemize}
\end{KeyConcept}

\begin{center}
\begin{tabularx}{\textwidth}{>{\raggedright\arraybackslash}p{2.6cm} >{\raggedright\arraybackslash}p{2.6cm} Y >{\raggedright\arraybackslash}p{2.9cm}}
\rowcolor{AccentDark}
\thead{Framework} & \thead{Focus} & \thead{Core question} & \thead{Key keyword} \\
\toprule
Deontology & Duty / means & Were the means right? & Right duty; pure means \\
Teleology / Consequentialism & Consequences / ends & Were the consequences right? & Good outcome; welfare maximisation \\
Virtue Ethics & Character & What kind of person am I becoming? & Virtues and character formation \\
\bottomrule
\end{tabularx}
\end{center}

\subsection{Applied (Practical) Dimension}

Applied ethics is the application of ethical principles to specific domains. In every particular context, certain ethical standards become more or less important. The applied dimension asks: given these theoretical frameworks, which principles should guide decision-making in this specific area?

Applied ethics includes: environmental ethics, technology ethics, media ethics, bioethics, sports ethics, professional ethics (medical, legal, engineering), and administrative ethics. Each is discussed in the sections below.

\section{Ethical Relativism and Ethical Absolutism}

\begin{center}
\begin{tabularx}{\textwidth}{>{\raggedright\arraybackslash}p{2.6cm} Y Y}
\rowcolor{AccentDark}
\thead{Basis} & \thead{Ethical Relativism} & \thead{Ethical Absolutism} \\
\toprule
Core claim & Ethical standards are contingent upon circumstances, time, and place & Ethical standards are absolute --- applicable universally, irrespective of circumstance, time, or place \\
Implication & Right and wrong vary across societies and across time & Certain standards are universally binding on all human beings \\
\bottomrule
\end{tabularx}
\end{center}

\subsection{Ethical Relativism}

Ethical relativism holds that ethical standards are not absolute but contextual --- they differ from society to society and from era to era. What is considered ethical in one time or place may be considered unethical in another.

\begin{ExampleBox}
\begin{itemize}
\item Eating beef is widely acceptable in Western societies but not ethical in the Indian context; McDonald's does not serve beef in India but does in the USA.
\item Dress norms vary --- clothes considered decent in one society may be regarded as indecent in another.
\item Polygamy is practised in certain societies as a culturally sanctioned norm.
\item Certain tribal communities do not wear clothing; judging them by mainstream urban standards would be imposing one culture's values on another.
\item EWS reservation did not exist initially; as ethical understandings of economic fairness evolved, it was introduced.
\item CSR was not legally mandatory before the Companies Act, 2013 --- but evolved into a mandatory ethical obligation.
\item The Juvenile Justice Act originally did not allow any juvenile below 18 to be tried as an adult even for heinous crimes; ethical standards evolved to allow it for juveniles aged 16--18 who have attained early maturity.
\item Naturist communities (those who prefer to live without clothing) operate on a different but internally coherent set of norms.
\end{itemize}
\end{ExampleBox}

\paragraph{Why Ethical Relativism Matters --- Avoiding Ethnocentrism and Cultural Imperialism} Ethnocentrism means judging another culture through the lens of one's own ethical standards --- treating one's own culture as the benchmark. The danger is imposing one culture's values on others, which is called cultural imperialism.

\begin{ExampleBox}
\begin{itemize}
\item Colonial powers justified their colonialism as the ``White Man's Burden'' --- the supposed duty to ``civilise'' the colonised peoples. This is classic cultural imperialism: imposing one's own ethical framework as universal and superior.
\item Anthropologist Margaret Mead's study of Samoan society --- she found that sexual norms were far more liberal there than in mainstream American society. The correct approach is to understand, not to impose one's own standards.
\end{itemize}
\end{ExampleBox}

Ethical relativism is the intellectual defence against both ethnocentrism and cultural imperialism. We should not judge others' ethical standards through the lens of what we consider right and wrong, and we should not impose our standards on them in the name of ethical superiority.

\paragraph{Standardisation Due to ICT and Globalisation} Information and communication technologies (ICT), the internet, and advocacy groups are increasingly standardising ethical norms across societies. Practices once considered culturally acceptable are being questioned in the light of evolved universal standards.

\begin{ExampleBox}
\begin{itemize}
\item Triple talaq --- once considered ethically acceptable in one community, now rendered unjust by law.
\item Sabarimala --- the exclusion of women from a sacred space was challenged and changed.
\item Sati Pratha --- abolished through social reform and legislation.
\item Shikari Utsav in West Bengal --- certain communities celebrate this traditional hunting festival; but in the evolved context of wildlife conservation, this practice is being questioned.
\item Opening of temple sacred spaces to all sections of society.
\end{itemize}
\end{ExampleBox}

\paragraph{Merits and Demerits of Ethical Relativism}

\begin{center}
\begin{tabularx}{\textwidth}{Y Y}
\rowcolor{AccentDark}
\thead{Merits} & \thead{Demerits} \\
\toprule
Promotes tolerance of cultural diversity & Can lead to moral nihilism if taken to an extreme --- if every standard is relative, then nothing can be criticised \\
Avoids ethnocentrism and cultural imperialism & May actually promote intolerance: e.g. Taliban-style practices justified as ``our culture'' \\
Allows ethical standards to evolve to suit changing needs of the time & Undermines moral progress --- social reformers could not have challenged Sati Pratha if all standards were relative \\
Respects context in applying universal principles (e.g. CBDR in climate) & There are some universal standards that cannot be relativised (dignity, justice, fairness) \\
\bottomrule
\end{tabularx}
\end{center}

\begin{CriticalPoint}
The biggest problem with extreme ethical relativism is moral nihilism: if everyone is right by their own standards (the person who says ``I am right, the world is dissatisfied with even God''), then there are no standards left at all. Taliban will say ``this is our culture''; Kasab may believe he was doing the right thing; a person who throws acid can claim justification in personal grievance. Once you allow each person or society to define ethics for themselves without any external standard, you cannot condemn anything. This is the fundamental demerit.
\end{CriticalPoint}

\subsection{Ethical Absolutism}

Ethical absolutism holds that certain ethical standards are absolute --- they apply universally, irrespective of cultural context, time, or place. Dignity, justice, and fairness are examples of such standards. Wherever human beings are, they should be accorded basic dignity.

\begin{ExampleBox}
Imagine two workers A and B in any society anywhere in the world: both are of equal ability, work equally long, produce equal output --- but one is paid X and the other 100X. In every society, the underpaid worker will experience a sense of unfairness. This universality of the sense of injustice is what ethical absolutists argue for --- fairness and justice are absolute standards.
\end{ExampleBox}

\subsection{Reconciling Relativism and Absolutism}

In practice, both absolute and relative ethical standards coexist --- just as, in the Indian Constitution, certain parts (the basic structure) cannot be changed while other parts can be.

\begin{KeyConcept}
\begin{itemize}
\item Universal principles (dignity, justice, fairness) are absolute and must be respected everywhere.
\item Application of those principles must be contextual and local.
\item ``Common But Differentiated Responsibility'' (CBDR) in climate change: sustainability is a universal principle, but its application differs --- developed countries have greater responsibility than developing countries because of their historical emissions.
\item Juvenile Justice Act: justice is an absolute principle, but its application to a juvenile must account for local context (early maturity, heinous crime) --- otherwise there will be a travesty of justice.
\end{itemize}
\end{KeyConcept}

\begin{PYQLink}
UPSC question (POCSO Act case): a 13-year-old girl was found to have ``married'' a 25-year-old man; under POCSO there is no concept of consensual sex when one party is a minor. The man was held guilty. However, an expert panel found the victim remained committed to the accused and faced hardship from police, society, and the legal system. The Supreme Court, exercising Article 142 powers, withheld sentencing and directed the West Bengal government to ensure her welfare --- because pronouncing punishment would create even greater hardship for the family. This illustrates: letter demands punishment; spirit and context demand a more sensitive application of justice. Universal principle (protect children) must be applied contextually (avoid compounding victim's suffering).
\end{PYQLink}

\begin{QuickRevision}
\begin{itemize}
\item Ethical relativism: standards are contingent on context, time, society. Merits: tolerance, diversity, evolution of norms. Demerits: moral nihilism, undermines moral progress.
\item Ethical absolutism: certain standards (dignity, justice, fairness) are universal and non-negotiable.
\item Best approach: believe in universal principles; apply them locally and contextually. CBDR is the model.
\end{itemize}
\end{QuickRevision}

\section{Environmental Ethics (Applied Ethics)}

Applied ethics is the application of ethical principles to specific domains. Environmental ethics applies them to our relationship with the natural environment. In every environmental issue, certain ethical principles become especially important and should guide decision-making.

\subsection{Shallow Ecology vs Deep Ecology}

\begin{center}
\begin{tabularx}{\textwidth}{>{\raggedright\arraybackslash}p{2.8cm} Y Y}
\rowcolor{AccentDark}
\thead{Basis} & \thead{Shallow Ecology} & \thead{Deep Ecology} \\
\toprule
Core view & Environment is important only insofar as it benefits humans; it has no intrinsic value of its own & Environment has intrinsic worth --- irrespective of its contribution to human welfare \\
Decision-making perspective & Anthropocentrism: human-centric decision-making & Ecocentrism: decision-making from the perspective of the entire ecosystem \\
Example & Solve air pollution because it harms our health & Protect a forest because the forest itself has worth, not just because we derive timber or tourism from it \\
\bottomrule
\end{tabularx}
\end{center}

\begin{CriticalPoint}
Anthropocentrism --- greed-driven, human-centric decision-making --- has driven climate change, declining soil fertility, falling groundwater levels, and deforestation. These are the consequences of always asking ``what's in it for us?'' without considering the entire ecosystem.
\end{CriticalPoint}

\subsection{Ecocentrism}

Ecocentrism holds that decision-making about the environment should be done not just from the perspective of human beings but from the perspective of the entire ecosystem --- all its biotic (humans, animals, plants) and abiotic (soil, rock, water) components and the relationships among them.

\begin{ExampleBox}
\begin{itemize}
\item Declaring a species as ``vermin'' (to be killed) is a decision that must be evaluated not just from the perspective of human convenience but from the perspective of the ecosystem as a whole. That species may play a critical role in the food chain or ecosystem stability.
\item AI-enabled cameras in Shimla for tiger reserves: detecting poaching automatically is an example of technology deployed from an ecocentric perspective --- protecting the ecosystem.
\item E-pass systems in Ooty and Kodaikanal regulating tourist activity in ecologically sensitive areas.
\end{itemize}
\end{ExampleBox}

\begin{PYQLink}
``Forests precede civilisations and deserts follow them.'' Forests provide food, fodder, timber and a host of ecosystem services that sustain civilisations. When humans destroy forests out of greed, they ultimately convert fertile land to desert. This is why ecocentrism, not anthropocentrism, must guide our decisions.
\end{PYQLink}

\subsection{Key Principles of Environmental Ethics}

\begin{itemize}
\item Stewardship: we hold the environment in trust. ``We do not inherit the environment from our ancestors; we borrow it from future generations.'' (Gandhiji.) As trustee of something borrowed, you take care of it, not destroy it. This is our obligation towards future generations.
\item Justice and equity: Polluter Pays Principle --- each entity that pollutes should pay the cost of the pollution. Common But Differentiated Responsibility (CBDR) --- all nations are responsible for environmental sustainability, but developed nations bear more responsibility because of their historical emissions. Justice to future generations is also demanded.
\item Precautionary Principle: when the consequences of a technology are not fully known, err on the side of caution. Do not deploy until the risks are assessed. This is why decisions about GM crops are slow --- the precautionary principle demands caution in the face of uncertainty.
\item Global perspective: environmental decisions have global significance. Stubble burning in Punjab creates smog in Delhi; US withdrawal from the Paris Climate Deal affects the entire world. We must transcend narrow perspectives and think globally.
\item Pacifism and peace: wars and conflicts cause massive environmental degradation --- missiles, explosions, displacement. Peace is a precondition for environmental protection.
\item Rights: animal rights are part of environmental ethics. Animals too have rights and should not be subjected to cruelty or extinction.
\item Duties --- 4R: Reduce, Reuse, Recycle, and Refuse (control consumerism). Citizens' duty to manage waste and minimise consumption.
\end{itemize}

\begin{GovtPolicy}
Extended Producer Responsibility (EPR): companies are responsible for taking back their electronic and plastic waste into the recycling system. This is an application of accountability and stewardship in environmental policy.
\end{GovtPolicy}

\begin{QuickRevision}
\begin{itemize}
\item Shallow ecology = anthropocentrism = human-centric. Deep ecology = ecocentrism = ecosystem-centric.
\item Key principles: stewardship, justice/equity (Polluter Pays, CBDR), precautionary principle, global perspective, pacifism, animal rights, 4R duties.
\item Apply these principles in all environmental case studies, section-A answers, and essays.
\end{itemize}
\end{QuickRevision}

\section{Technology Ethics (Applied Ethics)}

Technology ethics is the set of ethical principles that guide the development, deployment, and use of technology. Just as ``what are ethics?'' can be answered as ``standards that guide decision-making about right and wrong'', technology ethics = the standards that guide our decision-making in the context of technology.

\begin{CriticalPoint}
This is one of the most likely areas for a UPSC Mains question this year, especially on AI and deepfake. Prepare this section thoroughly.
\end{CriticalPoint}

\subsection{Key Principles of Technology Ethics}

\paragraph{Privacy} Technology, especially digital technology, poses grave threats to privacy. End-to-end encryption on messaging apps like WhatsApp is a privacy-protection mechanism. Online transactions must be secured. The GDPR in the European Union is the most significant privacy-protection framework globally, arising directly from the intrusive capability of modern technology.

\begin{ExampleBox}
Cybercriminals impersonating bank employees and extracting OTPs lead to accounts being emptied. Security protocols and informed consent about data use are privacy-protection tools.
\end{ExampleBox}

\paragraph{Transparency} Whenever technology is deployed to collect or use data, people must be told what that data will be used for. Data must only be used for the purpose for which it was collected. Any app or service that accesses your data must be transparent about how it will be used.

\paragraph{Accountability} Those who develop and deploy technology must be accountable for how it is used and misused. Technology companies are more powerful than individual users because they have proprietary knowledge of their own technology, creating an asymmetry that must be checked by accountability.

\begin{GovtPolicy}
\begin{itemize}
\item WhatsApp messaging limits: WhatsApp took responsibility by limiting forwards to 5 recipients to curb viral spread of fake news. This is an example of a technology company taking accountability for the social consequences of its platform.
\item Extended Producer Responsibility (EPR) in plastic and e-waste rules: companies must take back their products and recycle them --- accountability for the environmental consequences of their technology.
\item DiDi (Chinese ride-hailing app): fined \$600 million for sending user data to outside the country, violating transparency and data sovereignty norms.
\end{itemize}
\end{GovtPolicy}

\paragraph{Fairness --- Especially Critical for AI} AI algorithms are trained by feeding large datasets. If the training data is biased, the algorithm will be biased, and decisions made by that AI will perpetuate or amplify existing biases.

\begin{ExampleBox}
\begin{itemize}
\item PM Modi's comment at an AI meet: if you prompt an AI to generate an image of a person writing with the left hand, the AI was generating an image of a person writing with the right hand --- because the training data overwhelmingly contained right-handed writers. The algorithm was biased.
\item Amazon's AI hiring platform: the platform was trained on resumes of existing employees (who were predominantly male) and began systematically rejecting female candidates. The keyword ``cricket'' is more common in male resumes (since more males play cricket in India); the algorithm chose it as a positive indicator, further biasing towards men. Amazon shut the platform down.
\end{itemize}
\end{ExampleBox}

If AI is used in decision-making --- hiring, loan approval, bail decisions, medical diagnosis --- algorithmic bias can cause profound injustice. Fairness in AI requires diverse and representative training data, regular audits, and human oversight.

\paragraph{Beneficence} The primary intention of any technology should be to benefit humanity. Technology must be developed and deployed for the public good.

\begin{ExampleBox}
\begin{itemize}
\item AI-enabled CCTV cameras to monitor board examinations in Karnataka --- curbing malpractice.
\item AI-enabled cameras in Shimla tiger reserves --- automatically alerting rangers when poaching activity is detected.
\item ``Nudge'' letters sent by tax authorities to taxpayers: using data non-intrusively to nudge taxpayers into revising and updating their returns, generating \rupee 11,000 crore in tax revenue over four years. This is beneficent use of technology that guided and enabled rather than coerced.
\end{itemize}
\end{ExampleBox}

\paragraph{Non-Maleficence} ``If you cannot do good, at least do not do harm.'' When a technology developer is not sure of the consequences of their technology, they should refrain from deploying it, or deploy it cautiously under the precautionary principle.

\begin{ExampleBox}
Driverless cars in Indian road conditions: the unpredictability of Indian traffic, pedestrians, and infrastructure means that deploying driverless technology without adequate safeguards could cause serious harm. Non-maleficence demands caution.
\end{ExampleBox}

\paragraph{Informed Consent} Before any technology is used on people --- particularly in medical or data contexts --- people must be fully informed and must give their consent. This is what ``informed consent'' means: consent that is given with full knowledge of what is being agreed to.

\begin{ExampleBox}
Clinical trials for new drugs: participants must be fully informed about risks and must consent before participating --- informed consent in medical technology.
\end{ExampleBox}

\paragraph{Sustainability} Technology must be used sustainably, not in ways that deplete resources or damage the environment for short-term gain.

\begin{ExampleBox}
\begin{itemize}
\item Green Revolution --- groundwater pumps gave farmers access to irrigation, but farmers began growing rice in areas ecologically unsuited for it (Punjab), depleting groundwater at an unsustainable rate.
\item Indiscriminate use of urea (nitrogenous fertiliser) has disrupted N-P-K ratios in soil, reducing fertility over time. Sustainable technology use would have maintained the balance.
\item Sex-determination technology (ultrasound) --- intended for medical diagnosis but used for sex-selective abortions, leading to a skewed sex ratio. The technology is not inherently bad; the unsustainable and unethical use of it is the problem.
\end{itemize}
\end{ExampleBox}

\paragraph{Respect for Human Autonomy} Technology is a tool and should remain a servant of human beings, not become their master. Technology should support and expand human choices, not manipulate or override them.

\begin{KeyConcept}
``Science and technology is not a panacea for growth and security of the nation.'' Technology is only an instrument in human hands. Whether it creates problems or solves them depends on how humans use it: sustainably or unsustainably, transparently or opaquely, fairly or unfairly, accountably or unaccountably.
\end{KeyConcept}

\subsection{Challenges of Technology in Administration}

Technology is increasingly deployed to deliver public services and fight corruption. But there are significant ethical challenges in this reliance.

\begin{itemize}
\item Digital divide: 94\% of men vs.\ 74\% of women in Haryana have access to mobile phones (20\% gap). 63\% of elderly people do not know how to use the internet. If all services are delivered digitally, these groups are excluded.
\item Biometric failure: Aadhaar-based biometric authentication fails for manual labourers and the elderly whose fingerprints are worn or degraded. The PSE Committee flagged high failure rates of Aadhaar authentication. When biometrics fail, ration is denied --- people who are already vulnerable can go without food.
\item Algorithmic bias: as discussed, biased algorithms can entrench existing injustices.
\item Empathy deficit: AI-driven administration lacks the compassion and empathy of human decision-making, which can be critical in sensitive cases.
\item Accountability gap: when an AI makes a wrong decision, who is responsible --- the algorithm, the company, or the officer who relied on it?
\item Malfunction: technology can malfunction; over-reliance creates fragility.
\end{itemize}

\begin{PYQLink}
UPSC question theme: ``Use of technology in governance: positive and negative ethical dimensions.'' Positive: AI cameras in exam halls, nudge letters for tax compliance, GeM portal for transparent procurement. Negative: digital divide excluding vulnerable groups, biometric failures denying rations, algorithmic bias, empathy deficit.
\end{PYQLink}

\begin{QuickRevision}
\begin{itemize}
\item Technology ethics principles: privacy, transparency, accountability, fairness (esp. for AI), beneficence, non-maleficence, informed consent, sustainability, respect for autonomy.
\item Challenges in administration: digital divide, biometric failure, algorithmic bias, empathy deficit, accountability gap, malfunction.
\item Technology is an instrument --- outcome depends on how humans use it.
\end{itemize}
\end{QuickRevision}

\section{Media Ethics (Applied Ethics)}

Media ethics are the standards that media organisations and journalists should uphold in the collection, processing, and transmission of information, whether through print, electronic, or digital channels.

\subsection{Key Principles of Media Ethics}

\paragraph{Accuracy / Verification} Media must verify information before transmitting it. Sharing unverified information can cause panic, social unrest, or injustice.

\begin{ExampleBox}
\begin{itemize}
\item Poonam Pandey case: media published the fake news of her death without verification, which was later revealed to be a stunt. Though her goal was to raise awareness about cervical cancer, the initial reporting demonstrates the danger of unverified information.
\item Operation Sindoor: the government had to issue a press advisory asking media not to post about live movement of army troops. Common sense and national security demand that such restraint is the right thing to do, even though media might argue they have the freedom to report. Freedom of expression with reasonable restrictions.
\end{itemize}
\end{ExampleBox}

\paragraph{Freedom with Reasonable Restrictions} Media has freedom of speech and expression, but that freedom comes with reasonable restrictions --- particularly on grounds of national security, sovereignty, public order, and the dignity of individuals. You may have the right to publish something, but it may not be the right thing to do.

\begin{KeyConcept}
``You may have a right to do something, but it may not be the right thing to do.'' --- right vs right interface.
\end{KeyConcept}

\paragraph{Independence and Autonomy} Media acts as the fourth pillar of democracy. This role requires independence from government and corporate power. Media that is not independent cannot hold power accountable.

\begin{ExampleBox}
\begin{itemize}
\item Sedition law / UAPA used to threaten journalists: when the state uses legal action as a tool to silence criticism, it threatens media independence.
\item Government using print media advertisements as a reward (to compliant media) or withholding them as punishment (from critical media).
\item Corporate takeover of media houses: when a single corporate entity controls multiple media houses, the wrongdoings of that corporate will not be reported --- again threatening independence.
\end{itemize}
\end{ExampleBox}

\paragraph{Impartiality and Fairness} Media must be impartial in its coverage and fair to all those it reports on. When media makes up its mind before hearing all sides, it prejudges cases.

\begin{ExampleBox}
\begin{itemize}
\item Media trials in cases like Rhea Chakravarti or Arushi Darwar: media had already decided who the villain was before any court judgment. This is a gross failure of fairness and goes against the principle of presumption of innocence.
\item Right of reply: if you are writing critically about a minister or any person, you must give them the opportunity to present their own perspective. This is a requirement of fairness.
\end{itemize}
\end{ExampleBox}

\paragraph{Humanity / Compassion} Coverage should be conducted with basic human compassion. People in trauma or grief deserve personal space and dignity, even if they are public figures.

\begin{ExampleBox}
When Sushant Singh Rajput died, media arrived at his family's home and pushed microphones into his grieving father's face to ask ``How do you feel?'' People should be allowed to grieve in personal space. Coverage must be done with compassion, not just to chase TRP.
\end{ExampleBox}

\paragraph{Do Not Incite Hatred} Media must not inflame hatred on grounds of nationality, religion, caste, community, or gender.

\begin{ExampleBox}
Sudarshan TV --- UPSC Jihad programme: the channel ran a programme claiming to expose a ``conspiracy'' of Muslim candidates in the civil services. The Supreme Court found it to be communal propaganda. TRAI took action. This is a clear case of inciting hatred on grounds of religion under the guise of news.
\end{ExampleBox}

\paragraph{Public Interest over Sensationalism} Media should prioritise stories that matter to the public --- NEET irregularities, policy failures, corruption --- rather than sensationalist content or celebrity gossip. Clickbait thumbnails and TRP-driven content are not serving public interest.

Similarly, media must be able to distinguish facts from opinions. Presenting conjecture or personal interpretation as verified fact is a serious ethical violation. Social media amplifies this problem: anonymity allows people to say anything without accountability.

\begin{QuickRevision}
\begin{itemize}
\item Media ethics principles: accuracy/verification, freedom with reasonable restrictions, independence \& autonomy, impartiality/fairness, humanity/compassion, no incitement of hatred, public interest over sensationalism.
\item Distinguish: facts (verified, objective) vs. opinions (conjecture, interpretation).
\item Negative examples: media trials (Rhea Chakravarti), Sudarshan TV communal propaganda, unverified Poonam Pandey death story, reporter pushing mic at grieving father.
\end{itemize}
\end{QuickRevision}

\section{Ethics in Public and Private Relationships}

\subsection{Distinction Between Public and Private Relationship}

\begin{center}
\begin{tabularx}{\textwidth}{>{\raggedright\arraybackslash}p{2.9cm} Y Y}
\rowcolor{AccentDark}
\thead{Basis} & \thead{Private Relationship} & \thead{Public Relationship} \\
\toprule
Scope & Relationships with family members and friends & All other interactions --- as a citizen, professional, or public official (Delhi Metro, traffic, cricket, public office) \\
Essence & Love, human emotion, personal attachment & Societal/collective interest and public welfare \\
Primary source of ethical standards & Traditions, customs, family norms, cultural practices & Laws, rules, regulations, professional codes \\
Nature of enforcement & Informal, forgiving --- driven by love and emotion & Formal --- police, judiciary, professional bodies \\
Key ethical obligation & Loyalty, care, forgiveness, emotional support & Selflessness, impartiality, accountability, public interest \\
\bottomrule
\end{tabularx}
\end{center}

\subsection{Role of Traditions in Private Relationship Ethics}

Our understanding of how to behave towards parents, spouses, siblings, and friends is mostly shaped by family norms, customs, and traditions --- not by formal law.

\begin{ExampleBox}
\begin{itemize}
\item Marriage vows / Saptapadi (Hindu wedding) or the equivalent in other faiths --- these formalise the ethical obligations between spouses: how a husband should treat a wife and vice versa.
\item Rakshabandhan: formalises the obligation of siblings to protect and support each other.
\item Shraddha Paksha / ancestor rituals: obligations to pay respect to ancestors --- an ethical obligation within the family extended even beyond life.
\end{itemize}
\end{ExampleBox}

\subsection{Why Private Relationships Are More Informal}

Private relationships are governed by love and emotion. Because love is the basis, forgiveness is more readily available; mistakes are tolerated; the informal standard is high but the enforcement is soft. By contrast, public relationships are governed by collective interest --- and the consequence of ethical failures in public life can affect many people, sometimes fatally.

\begin{ExampleBox}
A civil servant watching a TV series during office hours is the same as a lazy student. But the consequence is not limited to the officer alone --- the people dependent on that public service suffer. A soldier must prioritise duty even at the cost of personal pleasure or personal safety --- because lives depend on it.
\end{ExampleBox}

\begin{CriticalPoint}
In public relationships, the ethical obligation is to dissolve individual interest into public interest. Public servants are expected to follow a higher set of ethical standards precisely because the consequences of their conduct affect large numbers of people who are dependent on them.
\end{CriticalPoint}

\subsection{Consequences in the Context of Public vs Private}

When a student fails to study, the immediate consequence is personal. But when a public servant fails to do their duty, it compounds into systemic harm --- delayed projects, denied services, lives lost. This is why the ethical standard expected in public office is significantly higher than in private life, and why public servants are required to be selfless even at personal cost.

\begin{QuickRevision}
\begin{itemize}
\item Private relationship ethics: governed by traditions, family norms, customs; enforcement is informal and driven by love.
\item Public relationship ethics: governed by law, rules, professional codes; enforcement is formal.
\item Public office demands selflessness, impartiality, and prioritisation of public interest above personal interest.
\item The higher the consequence of ethical failure for others, the higher the ethical standard required of the actor.
\end{itemize}
\end{QuickRevision}
